\clearpage\begingroup\NoSourceHead
\makeatletter\tagsleft@true\let\veqno\@@leqno\makeatother
\renewcommand{\thefootnote}{\arabic{footnote})}\setcounter{footnote}{18}
\ArithmeticHead{14}{L}{BOLZA, GAUSS UND DIE VARIATIONSRECHNUNG.}
\noindent\hspace*{\parindent}7. Der erste Versuch, die Methode der Variation der unabhängigen Variabeln
auf zweidimensionale Probleme anzuwenden, rührt von \textsc{Euler}\footnote{\textsc{Euler}, loc. cit. 2), Art. 141—149.} (1770)
her. Er deutet den Variationsprozess geometrisch: den Variabeln $x,y,z$, die
als rechtwinklige Koordinaten eines Punktes einer Fläche $z=z(x,y)$ aufgefasst
werden, werden unendlich kleine Zuwächse $\delta x$, $\delta y$, $\delta z$ erteilt, die Funktionen
der beiden unabhängigen Variabeln $x,y$ sind. Dadurch entsteht eine der
ersten unendlich benachbarte Fläche, wobei zugleich dem Punkt $x,y,z$ der
ursprünglichen Fläche der Punkt $x+\delta x$, $y+\delta y$, $z+\delta z$ der zweiten Fläche
zugeordnet ist.

Es kommt nun zunächst auf die Berechnung der Variation der partiellen
Ableitungen
\[
p=\frac{\partial z}{\partial x},\quad q=\frac{\partial z}{\partial y}
\]
an. \textsc{Euler} schliesst aus der Bedeutung von $\delta p$, dass
\[
\delta p=\frac{\partial(z+\delta z)}{\partial(x+\delta x)}-\frac{\partial z}{\partial x}
=\frac{\partial z+\partial\delta z}{\partial x+\partial\delta x}-\frac{\partial z}{\partial x}
\]
und erhält daraus durch Vernachlässigung unendlich kleiner Grössen höherer
Ordnung
\begin{equation*}
\delta p=\frac{\partial\delta z}{\partial x}-p\frac{\partial\delta x}{\partial x}\tag{14}
\end{equation*}
und die analoge Formel für $\delta q$.

Das ist dasselbe Resultat, dass sich auch durch rein mechanische Anwendung
des $\delta$-Algorithmus nach der Formel
\[
\delta\left(\frac{\partial z}{\partial x}\right)
=\frac{\partial x\,\delta\partial z-\partial z\,\delta\partial x}{\partial x^2}
\]
unter Benutzung der Gleichungen $\delta\partial x=\partial\delta x$, $\delta\partial z=\partial\delta z$ ergeben haben würde,
\par\clearpage\SourceLeaf{472}
\ArithmeticHead{15}{R}{GESCHICHTE DER EXTREMA VON DOPPELINTEGRALEN VOR GAUSS.}
\noindent und in der Tat haben später \textsc{Lagrange}\footnote{Oeuvres de \textsc{Lagrange}, Bd. XI, S. 105: »On aura, en différentiant« (nämlich mit dem Symbol $\delta$).} und \textsc{Lacroix}\footnote{\textsc{Lacroix}, loc. cit. 18), Art. 861: »différentiant ces fractions à la manière ordinaire, et changeant un $d$ en $\delta$«.} die Formel (14) auf
diese Weise abgeleitet.

Aus den Formeln für $\delta p$, $\delta q$ leitet dann \textsc{Euler} die Variationen der zweiten
partiellen Ableitungen her, indem er diese schreibt
\[
r=\frac{\partial p}{\partial x},\quad s=\frac{\partial p}{\partial y}=\frac{\partial q}{\partial x},\quad t=\frac{\partial q}{\partial y}.
\]

Die Variation von $\delta s$ kann man dabei auf zwei verschiedenen Wegen
erhalten, die jedoch zu verschiedenen Resultaten führen. Um den Widerspruch
aufzuklären, sucht nun \textsc{Euler} nachzuweisen, dass \textls[100]{$\delta x$ nur von $x$, $\delta y$ nur von
$y$ abhängen dürfen}. Bei dieser spezialisierenden Annahme fallen dann die
beiden verschiedenen Werte von $\delta s$ zusammen.

Aber \textsc{Euler} traut der Sache selbst nicht recht; er schreibt: »Omnibus
autem dubiis in hac investigatione felicissime occurremus, si soli quantitati $z$
variationes tribuamus, binis reliquis $x$ et $y$ plane invariatis relictis, ita ut sit
tam $\delta x=0$ quam $\delta y=0$«. Er gibt daher, um sicher zu gehen, bei der
weiteren Behandlung der Doppelintegrale, über die wir bereits in Nr. 1 berichtet
haben, die Variation der unabhängigen Variabeln $x,y$ auf und variiert
nur die Funktion $z$; und dabei ist er auch in seiner späteren Arbeit\footnote{\textsc{Euler}, loc. cit. 3).} vom
Jahr 1771 geblieben. Damit hat er sich freilich, da er auch die Grenzen
nicht variiert, den Zugang zu Problemen mit variablem Integrationsbereich
versperrt, für die übrigens die Zeit ohnehin nicht reif war. —

Die Überlegungen, durch die \textsc{Euler} zu beweisen sucht, dass man $\delta x$ als
nur von $x$, $\delta y$ nur von $y$ abhängig annehmen müsse, sind nicht stichhaltig,
und der wahre Grund des Widerspruchs der beiden Ausdrücke für $\delta s$ liegt
darin, dass schon die \textsc{Euler}schen Formeln für $\delta p$, $\delta q$ falsch sind. Den Fehler
in \textsc{Euler}s Schlussweise hat erst \textsc{M. Ostrogradsky}\footnote{\textit{Mémoire sur le calcul des variations des intégrales multiples}, Journal für r. u. a. Math., Bd. 15 (1836), S. 334.} (1834) aufgedeckt.

Zugleich bestätigt sich hier, dass die oben als »formale Variation« bezeichnete
rein mechanische Anwendung des $\delta$-Algorithmus bei gleichzeitiger
Variation der unabhängigen Variabeln, unter Umständen zu direkt falschen
\par\clearpage\SourceLeaf{473}
\ArithmeticHead{16}{L}{BOLZA, GAUSS UND DIE VARIATIONSRECHNUNG.}
\noindent Resultaten führen kann, und daher eine ohne genaue Feststellung der Grenzen
ihrer Anwendbarkeit unzulässige Operation ist.

8. Nach diesem ersten missglückten Versuch von \textsc{Euler} ist erst vierzig
Jahre später das Problem der Variation von Doppelintegralen bei variablem
Integrationsbereich wieder aufgenommen worden und zwar von \textsc{Lagrange}\footnote{Oeuvres de \textsc{Lagrange}, Bd. XI, S. 103—110.}
selbst in einem kurzen Exkurs über Variationsrechnung in der zweiten Auflage
der \textit{Mécanique analytique} (1811).

Es handelt sich darum, die erste Variation des Doppelintegrals
\[
J=\iint V(x,y,z,p,q,\ldots)\,dx\,dy
\]
bei gleichzeitiger Variation von $x,y,z$ zu berechnen. Dabei macht \textsc{Lagrange}
dieselbe spezialisierende Annahme wie \textsc{Euler}, dass $\delta x$ nur von $x$, $\delta y$ nur von
$y$ abhängt, dass somit
\begin{equation*}
\frac{\partial\delta x}{\partial y}=0,\quad\frac{\partial\delta y}{\partial x}=0,\tag{15}
\end{equation*}
ohne jedoch zu behaupten, dass diese Einschränkung notwendig sei, sondern
nur »um die Rechnung zu vereinfachen«. Es kommt nun vor allem wieder
auf die Berechnung von $\delta p$, $\delta q$ an. Nachdem \textsc{Lagrange}, wie bereits oben
erwähnt, die \textsc{Euler}sche Formel (14) für $\delta p$ durch »formale Variation« abgeleitet
hat, formt er dieselbe dadurch um, dass er das unter der speziellen Voraussetzung
(15) verschwindende Glied
\[
-q\frac{\partial\delta y}{\partial x}
\]
hinzufügt. So erhält er für $\delta p$, $\delta q$ die Ausdrücke
\begin{equation*}
\left\{\begin{aligned}
\delta p&=\frac{\partial\omega}{\partial x}+\frac{\partial p}{\partial x}\delta x+\frac{\partial q}{\partial x}\delta y,\\
\delta q&=\frac{\partial\omega}{\partial y}+\frac{\partial p}{\partial y}\delta x+\frac{\partial q}{\partial y}\delta y,
\end{aligned}\right.\tag{16}
\end{equation*}
wo nunmehr
\begin{equation*}
\omega=\delta z-p\delta x-q\delta y.\tag{17}
\end{equation*}

Aus den Formeln für $\delta p$, $\delta q$ leitet er dann sehr einfache und übersichtliche
Ausdrücke für die Variationen der höheren Ableitungen von $z$ her, mit deren
\par\clearpage\SourceLeaf{474}
\ArithmeticHead{17}{R}{GESCHICHTE DER EXTREMA VON DOPPELINTEGRALEN VOR GAUSS.}
\noindent Hilfe er nunmehr $\delta V$ auf die einfache Form bringen kann
\begin{equation*}
\delta V=\frac{\partial V}{\partial x}\delta x+\frac{\partial V}{\partial y}\delta y+N\omega
+P\frac{\partial\omega}{\partial x}+Q\frac{\partial\omega}{\partial y}+\ldots.\tag{18}
\end{equation*}

Weiter gibt \textsc{Lagrange} ohne weitere Erläuterung, also doch wohl auf Grund
»formaler Variation«, für die Variation des Integrals $J$ den Ausdruck
\begin{equation*}
\delta J=\iint\delta(V\,dx\,dy)=\iint(dx\,dy\,\delta V+V\delta(dx\,dy)).\tag{19}
\end{equation*}

Die nächste Schwierigkeit ist die Berechnung von $\delta(dx\,dy)$. Hier geht
\textsc{Lagrange} auf die geometrische Bedeutung von $dx\,dy$ als Flächeninhalt des
Rechtecks mit den Eckenkoordinaten
\[
x,\ y;\quad x+dx,\ y;\quad x,\ y+dy;\quad x+dx,\ y+dy
\]
zurück und erhält
\begin{equation*}
\delta(dx\,dy)=\left(\frac{\partial\delta x}{\partial x}+\frac{\partial\delta y}{\partial y}\right)dx\,dy.\tag{20}
\end{equation*}

Die Formel (19) geht nunmehr über in
\begin{equation*}
\begin{aligned}
\delta J={}&\iint\left(\delta V-\frac{\partial V}{\partial x}\delta x-\frac{\partial V}{\partial y}\delta y\right)dx\,dy\\
&+\iint\left(\frac{\partial(V\delta x)}{\partial x}+\frac{\partial(V\delta y)}{\partial y}\right)dx\,dy;
\end{aligned}\tag{21}
\end{equation*}
dabei sind, — ebenso wie auch schon in (18) —, vor Ausführung der Differentiationen
nach $x$ und $y$ die Grössen $z,p,q,\ldots$ als Funktionen von $x,y$ in
$V$ eingesetzt zu denken.

Diese Formel ist für die Methode der Variation der unabhängigen Variabeln
das Gegenstück zu der Fundamentalformel (9) von \textsc{Ohm}; wir werden
sie die \textsc{Lagrange}sche \textls[100]{Fundamentalformel} nennen. Sie ist später von
\textsc{Ostrogradsky}\footnote{\textsc{Ostrogradsky}, loc. cit. 23), S. 339, 340.} (1834) und \textsc{G. M. Pagani}\footnote{\textit{Résolution d'un problème relatif au calcul des variations}, Journal für r. u. a. Mathematik, Bd. 15 (1836), S. 84 ff., insbes. S. 92.} (1834) auf $n$-fache Integrale
verallgemeinert worden.

Das Einsetzen des Wertes (18) von $\delta V$ liefert nun unter Benutzung der
Formeln (2) für $\delta J$ das \textls[100]{Schlussresultat}\footnote{\textsc{Lagrange} schreibt übrigens dieses Schlussresultat gar nicht explizite hin, ebensowenig wie die explizite Formel (21); er begnügt sich, den Gang der dazu führenden Rechnung angegeben zu haben.}:
\begin{equation*}
\delta J=\iint\Omega\omega\,dx\,dy+\iint\left(\frac{\partial(A+V\delta x)}{\partial x}+\frac{\partial(B+V\delta y)}{\partial y}\right)dx\,dy,\tag{22}
\end{equation*}
\par\smallskip\noindent{\small X 2 Abh. 5.}\hfill{\small 3}\par
\clearpage\SourceLeaf{475}
\ArithmeticHead{18}{L}{BOLZA, GAUSS UND DIE VARIATIONSRECHNUNG.}
\noindent wobei die Grössen $\Omega,A,B,\omega$ durch die Gleichungen (4) und (17) definiert
sind.

Die Formel (22) ist das genaue Analogon einer entsprechenden, von \textsc{Euler}\footnote{\textsc{Euler}, loc. cit. 2), Art. 86.}
für einfache Integrale bewiesenen Formel, die vermutlich \textsc{Lagrange} bei seiner
Untersuchung, insbesondere bei seiner Ableitung der Formeln (16) als Wegweiser
gedient hat.

Auf die weitere Behandlung des zweiten Doppelintegrals in (22) ist \textsc{Lagrange}
nicht eingegangen; es genügt ihm, festzustellen, dass dieses bei fester
Begrenzung verschwindet, da es ihm bei der ganzen Ableitung lediglich darauf
ankommt zu beweisen, dass man bei gleichzeitiger Variation von $x,y,z$ zu
derselben partiellen Differentialgleichung kommt, wie bei alleiniger Variation
von $z$.
\par\clearpage\endgroup
